\documentclass[12pt,a4paper]{article}
\usepackage[utf8]{inputenc}
\usepackage[T1]{fontenc}
\usepackage[polish]{babel}
\usepackage{geometry}
\usepackage{enumitem}
\usepackage{titlesec}
\usepackage{parskip}
\usepackage{fancyhdr}
\usepackage{xcolor}
\usepackage{mdframed}
\usepackage{microtype}
\usepackage{booktabs}
\usepackage{hyperref}
\usepackage{multirow}
\usepackage{array}
\usepackage{longtable}
\usepackage{tabularx}

\geometry{margin=2.5cm}
\hypersetup{colorlinks=true, linkcolor=black, urlcolor=blue}

\pagestyle{fancy}
\fancyhf{}
\rhead{SDLC-LLM Benchmark}
\lhead{Zadanie R3 -- Conflict detection}
\rfoot{\thepage}

\titleformat{\section}{\large\bfseries}{}{0em}{}[\titlerule]
\titleformat{\subsection}{\normalsize\bfseries}{}{0em}{}
\titleformat{\subsubsection}{\normalsize\itshape}{}{0em}{}

\definecolor{metabg}{gray}{0.93}
\definecolor{promptbg}{RGB}{235,245,255}
\definecolor{score3}{RGB}{220,240,220}
\definecolor{score2}{RGB}{255,243,205}
\definecolor{score1}{RGB}{255,220,200}
\definecolor{score0}{RGB}{255,200,200}
\definecolor{tpcolor}{RGB}{220,240,220}
\definecolor{fpcolor}{RGB}{255,220,200}
\definecolor{fncolor}{RGB}{255,243,205}

\newcommand{\scorebox}[1]{%
  \ifnum#1=3\colorbox{score3}{\textbf{3}}%
  \else\ifnum#1=2\colorbox{score2}{\textbf{2}}%
  \else\ifnum#1=1\colorbox{score1}{\textbf{1}}%
  \else\colorbox{score0}{\textbf{0}}\fi\fi\fi
}
\newcommand{\tp}[1]{\colorbox{tpcolor}{\small #1}}
\newcommand{\fp}[1]{\colorbox{fpcolor}{\small #1}}
\newcommand{\fn}[1]{\colorbox{fncolor}{\small #1}}

\begin{document}

% ─── STRONA TYTUŁOWA ────────────────────────────────────────────────
\begin{center}
  {\LARGE\bfseries Requirements Artifact}\\[1.5em]
  {\large \textbf{Case study:} Appointment Booking System}\\[0.3em]
  {\large \textbf{Model:} \texttt{claude-sonnet-4-6}}\\[0.3em]
  {\large \textbf{Tryb:} LLM-only}\\[0.3em]
  {\large \textbf{Wejście:} Surowy opis (challenging -- celowe konflikty)}
\end{center}

\vspace{1em}\hrule\vspace{1.5em}

% ─── PROMPT ─────────────────────────────────────────────────────────
\section{Prompt}

\begin{mdframed}[backgroundcolor=promptbg, linewidth=0pt, innerleftmargin=10pt, innerrightmargin=10pt, innertopmargin=8pt, innerbottommargin=8pt]
\textbf{User prompt:}\\[0.3em]
\small
Przeanalizuj poniższy surowy opis systemu:

<opis>
System umożliwia użytkownikom rezerwację wizyt u specjalistów (np. lekarzy lub konsultantów). Użytkownicy mogą przeglądać dostępne terminy, dokonywać rezerwacji oraz je anulować. Specjaliści zarządzają swoim grafikiem. System powinien zapobiegać konfliktom rezerwacji, ale w niektórych przypadkach dopuszczalne są wyjątki. Użytkownik może anulować wizytę w uzasadnionych przypadkach, jednak anulowanie może być ograniczone czasowo. Specjalista może modyfikować swój grafik, także dla już istniejących rezerwacji. System powinien obsługiwać różne role użytkowników i być skalowalny.
</opis>

Zwróć wynik w formacie Markdown, zachowując dokładnie dwie poniższe sekcje:
### 1. Lista wykrytych konfliktów
Wypunktuj wszystkie znalezione konflikty w dostarczonym opisie. Przy każdym punkcie krótko (jedno zdanie) uzasadnij, na czym polega problem.
### 2. Poprawiona wersja
Sformułuj nowy, spójny opis systemu w formie tekstu ciągłego, z którego usunięto wyżej wymienione konflikty, zachowując przy tym pierwotny cel biznesowy.

WYMAGANIA FORMATOWANIA (OBOWIĄZKOWE):
- W sekcji 1 każdy punkt MUSI zaczynać się od "- " (myślnik + spacja).
- W sekcji 1 NIE używaj numeracji (1., 2., itd.) ani innych symboli list (np. *, •).
- Sekcja 2 MUSI być tekstem ciągłym. NIE używaj w niej żadnych wypunktowań.
- Nie dodawaj żadnego tekstu wprowadzającego, podsumowań ani komentarzy poza dwiema wymaganymi sekcjami.
\end{mdframed}

% ─── WYKRYTE KONFLIKTY ──────────────────────────────────────────────
\section{Lista wykrytych konfliktów}

\begin{itemize}[leftmargin=1.5em, label={-}]
  \item System deklaruje zapobieganie konfliktom rezerwacji, jednocześnie dopuszczając nieokreślone „wyjątki”, co tworzy sprzeczność – nie wiadomo, kiedy ochrona przed konfliktem obowiązuje, a kiedy nie.
  \item Użytkownik może anulować wizytę „w uzasadnionych przypadkach”, ale jednocześnie anulowanie jest „ograniczone czasowo” – oba warunki mogą wzajemnie się wykluczać (np. uzasadniony przypadek pojawia się po upływie dozwolonego czasu).
  \item Specjalista może modyfikować grafik dla już istniejących rezerwacji, co stoi w sprzeczności z celem systemu polegającym na zapobieganiu konfliktom rezerwacji – modyfikacja może jednostronnie unieważnić potwierdzoną rezerwację użytkownika bez jego zgody.
  \item Brak definicji „uzasadnionych przypadków” dla anulowania przez użytkownika powoduje konflikt z regułą ograniczenia czasowego – system nie może jednocześnie egzekwować limitu czasowego i honorować nieokreślonych wyjątków od tej reguły.
  \item Specjalista może modyfikować grafik bez wskazania, co dzieje się z kolidującymi rezerwacjami, co uniemożliwia ukończenie procesu rezerwacji w sposób spójny.
\end{itemize}

% ─── POPRAWIONA WERSJA ──────────────────────────────────────────────
\section{Poprawiona wersja}

System umożliwia użytkownikom rezerwację wizyt u specjalistów, takich jak lekarze lub konsultanci. Użytkownicy mogą przeglądać dostępne terminy, dokonywać rezerwacji oraz je anulować wyłącznie do określonego momentu przed planowaną wizytą, zdefiniowanego jako nieprzekraczalny limit czasowy anulowania. Po upływie tego limitu anulowanie przez użytkownika nie jest możliwe, niezależnie od okoliczności. System bezwzględnie zapobiega konfliktom rezerwacji, nie dopuszczając podwójnych rezerwacji tego samego terminu u danego specjalisty. Specjaliści zarządzają swoim grafikiem i mogą modyfikować dostępne terminy, jednak zmiana terminu, który jest już objęty potwierdzoną rezerwacją, jest możliwa wyłącznie po uprzednim powiadomieniu użytkownika i uzyskaniu jego akceptacji lub po anulowaniu danej rezerwacji przez specjalistę z jednoczesnym poinformowaniem użytkownika o konieczności dokonania nowej rezerwacji. System obsługuje różne role użytkowników, w tym co najmniej rolę pacjenta/klienta oraz rolę specjalisty, i jest zaprojektowany z myślą o skalowalności, umożliwiającej obsługę rosnącej liczby użytkowników i specjalistów.

\newpage

% ─── METRYKI VS GOLD ────────────────────────────────────────────────
\section{Metryki porównawcze względem zdefiniowanych konfliktów}

\begin{mdframed}[backgroundcolor=metabg, linewidth=0pt, innerleftmargin=10pt, innerrightmargin=10pt, innertopmargin=8pt, innerbottommargin=8pt]
\small
\textbf{Legenda:} \tp{TP} = trafnie wykryty konflikt \quad \fp{FP} = halucynacja (wskazanie konfliktu, którego nie ma) \quad \fn{FN} = przeoczenie celowo wprowadzonego problemu\\[0.3em]
Tabela weryfikuje zdolność modelu do wykrycia 6 celowo wprowadzonych problemów w tekście źródłowym.
\end{mdframed}

\vspace{0.8em}

\begin{tabular}{p{7cm} p{6cm} l}
\toprule
\textbf{Celowo wprowadzony problem (Gold)} & \textbf{Identyfikacja modelu} & \textbf{Status} \\
\midrule
Brak definicji „uzasadnionych przypadków” & Wykryto (punkt 4) & \tp{TP} \\
Brak jednoznacznej reguły anulowania & Wykryto (punkt 2) & \tp{TP} \\
Konflikt: zapobieganie vs dopuszczanie konfliktów & Wykryto (punkt 1) & \tp{TP} \\
Brak definicji zachowania przy zmianie grafiku & Wykryto (punkty 3 i 5) & \tp{TP} \\
\midrule
Brak reguł overbooking & Przeoczono & \fn{FN} \\
Brak informacji o równoczesnych operacjach & Przeoczono & \fn{FN} \\
\bottomrule
\end{tabular}

\vspace{1em}
\begin{tabular}{lrrrr}
\toprule
& \textbf{TP} & \textbf{FP} & \textbf{FN} & \\
\midrule
Wykryte konflikty & 4 & 0 & 2 & \\
\midrule
\textbf{Precision} & \multicolumn{4}{l}{$4 / (4+0) = \mathbf{1{,}00}$} \\
\textbf{Recall}    & \multicolumn{4}{l}{$4 / (4+2) = \mathbf{0{,}67}$} \\
\textbf{F1 Score}  & \multicolumn{4}{l}{$2 \times 1{,}00 \times 0{,}67 / (1{,}00+0{,}67) = \mathbf{0{,}80}$} \\
\bottomrule
\end{tabular}

\vspace{0.5em}
\begin{mdframed}[backgroundcolor=metabg, linewidth=0pt, innerleftmargin=10pt, innerrightmargin=10pt, innertopmargin=6pt, innerbottommargin=6pt]
\small\textbf{Komentarz do analizy:} Model perfekcyjnie zidentyfikował wszystkie luki w logice biznesowej, stąd maksymalna precyzja (Precision = 1.00). Przeoczył jednak czysto inżynieryjne aspekty opisu, takie jak brak zabezpieczeń dla równoczesnych rezerwacji (concurrency) oraz reguł overbookingu, co obniżyło metrykę Recall do 0.67.
\end{mdframed}

% ─── METADANE I OCENA ───────────────────────────────────────────────
\section{Metadane i ocena}
\begin{table}[h]
\begin{tabularx}{\textwidth}{l X}
\toprule
\textbf{Pole} & \textbf{Wartość} \\
\midrule
Model & \texttt{claude-sonnet-4-6} \\
Tryb & LLM-only \\
Iteracja & 1 \\
Czas łącznie & 15,33s \\
\midrule
\multicolumn{2}{l}{\textit{Ocena (0--3)}} \\
\midrule
Correctness (M1)     & \scorebox{3} \quad wykryte konflikty są absolutnie trafne i świetnie uzasadnione \\
Completeness (M2)    & \scorebox{2} \quad pominięto inżynieryjne przypadki brzegowe (współbieżność/overbooking) \\
Consistency (M3)     & \scorebox{3} \quad "Poprawiona wersja" idealnie łata wszystkie znalezione przez model luki \\
Clarity (M4)         & \scorebox{3} \quad rygorystyczne przestrzeganie narzuconego formatowania \\
Maintainability (M5) & \scorebox{3} \quad zrewidowany tekst nadaje się bezpośrednio do analizy technicznej \\
\bottomrule
\end{tabularx}
\end{table}

\end{document}